\documentclass[11pt,letterpaper]{article}
\usepackage[utf8]{inputenc}
\usepackage[T1]{fontenc}
\usepackage[spanish,es-noshorthands,es-tabla]{babel}
\usepackage[margin=2cm,top=2.5cm,bottom=2.5cm]{geometry}
\usepackage{amsmath,amssymb}
\usepackage{enumitem}
\usepackage{fancyhdr}
\usepackage{listings}
\usepackage{xcolor}
\usepackage{tcolorbox}
\usepackage{booktabs}
\usepackage{array}
\raggedbottom
\setlength{\headheight}{14pt}
\let\vspaceoriginal\vspace
\renewcommand{\vspace}[1]{\par\vspaceoriginal{#1}\par}
\pagestyle{fancy}
\fancyhf{}
\fancyhead[L]{Basic C Without Tears}
\fancyhead[C]{\NombreDocumento}
\fancyhead[R]{Stack}
\fancyfoot[L]{Basic-C-Without-Tears}
\fancyfoot[R]{\thepage}
\renewcommand{\headrulewidth}{0.4pt}
\renewcommand{\footrulewidth}{0.4pt}
\lstdefinestyle{customc}{
  language=C,
  numbers=left,
  stepnumber=1,
  numbersep=8pt,
  numberstyle=\tiny\color{gray},
  basicstyle=\footnotesize\ttfamily,
  keywordstyle=\bfseries\color{blue!80!black},
  commentstyle=\itshape\color{green!50!black},
  stringstyle=\color{red!70!black},
  breaklines=true,
  breakatwhitespace=true,
  tabsize=4,
  showstringspaces=false,
  frame=none
}
\newtcolorbox{solucionbox}{
  before=\par,
  after=\par,
  colback=white,
  colframe=black,
  arc=0mm,
  boxrule=0.6pt,
  left=3mm,
  right=3mm,
  top=2mm,
  bottom=2mm
}
\newif\ifpauta
\ifdefined\PAUTA\pautatrue\else\pautafalse\fi

\newcommand{\NombreDocumento}{Examen 02 -- Variante B}
\begin{document}
\begin{center}
    {\LARGE \textbf{Basic C Without Tears}} \\[0.2cm]
    {\Large \textbf{Examen 02 -- Variante B}} \\[0.12cm]
    \small Fundamentos de C, validación, matrices y cadenas de caracteres \\[0.08cm]
    \textbf{Autor:} Stack
\end{center}
\vspace{0.2cm}
\begin{enumerate}[label=\textbf{\arabic*.}]
\item Una biblioteca registra préstamos de libros en una matriz donde cada fila es una semana y cada columna una sucursal:
\begin{lstlisting}[style=customc, numbers=none]
int P[2][3] = {{4, 7, 3}, {5, 2, 8}};
\end{lstlisting}
\begin{enumerate}[label=(\alph*)]
\item Indique el valor de \texttt{P[0][2]} y \texttt{P[1][1]}.
\ifpauta\begin{solucionbox}\texttt{P[0][2]} vale 3 y \texttt{P[1][1]} vale 2.\end{solucionbox}\else\vspace{1.5cm}\fi
\item Calcule el total de préstamos de cada sucursal (cada columna).
\ifpauta\begin{solucionbox}Sucursal 1: $4+5=9$; sucursal 2: $7+2=9$; sucursal 3: $3+8=11$.\end{solucionbox}\else\vspace{2cm}\fi
\item Escriba un fragmento en C que recorra la matriz e imprima el total de cada columna.
\ifpauta\begin{solucionbox}\begin{lstlisting}[style=customc]
for (int j = 0; j < 3; j++) {
    int suma = 0;
    for (int i = 0; i < 2; i++) {
        suma += P[i][j];
    }
    printf("Sucursal %d: %d\n", j + 1, suma);
}
\end{lstlisting}Se fija una columna en el ciclo exterior y se recorren sus filas en el interior.\end{solucionbox}\else\vspace{3cm}\fi
\end{enumerate}
\newpage
\item Para ingresar una nota se usa la variable \texttt{float nota}; el rango permitido es de 1.0 a 7.0, inclusive.
\begin{enumerate}[label=(\alph*)]
\item Escriba la condición que detecta una nota fuera del rango.
\ifpauta\begin{solucionbox}\texttt{nota < 1.0f || nota > 7.0f}. Basta con que se cumpla una de las dos condiciones para que sea inválida.\end{solucionbox}\else\vspace{1.5cm}\fi
\item Escriba el fragmento \texttt{if-else} que imprima si la nota es válida.
\ifpauta\begin{solucionbox}\begin{lstlisting}[style=customc]
if (nota >= 1.0f && nota <= 7.0f) {
    printf("Nota valida\n");
} else {
    printf("Nota invalida\n");
}
\end{lstlisting}\end{solucionbox}\else\vspace{2cm}\fi
\item Clasifique las entradas \texttt{0.9}, \texttt{1.0}, \texttt{5.5} y \texttt{7.1}.
\ifpauta\begin{solucionbox}\texttt{0.9}: inválida; \texttt{1.0}: válida; \texttt{5.5}: válida; \texttt{7.1}: inválida.\end{solucionbox}\else\vspace{2cm}\fi
\end{enumerate}
\newpage
\item Determine qué imprime el programa.
\begin{lstlisting}[style=customc]
int A[2][2] = {{2, 4}, {6, 8}};
int suma = 0;
for (int i = 0; i < 2; i++) {
    for (int j = 0; j < 2; j++) {
        suma += A[i][j];
        printf("%d ", suma);
    }
}
printf("\nTotal=%d\n", suma);
\end{lstlisting}
\ifpauta\begin{solucionbox}El acumulado se imprime después de agregar cada elemento: 2, 6, 12 y 20.\par\textbf{Salida:}\\\texttt{2 6 12 20}\\\texttt{Total=20}\end{solucionbox}\else\vspace{3cm}\fi
\item Escriba una función que cuente cuántas cifras numéricas (\texttt{'0'} a \texttt{'9'}) contiene una cadena. No use funciones de biblioteca para strings.
\begin{center}\begin{tabular}{ll}\toprule\textbf{Entrada} & \textbf{Resultado}\\\midrule\texttt{"Aula 204"} & \texttt{3}\\\texttt{"C"} & \texttt{0}\\\texttt{"12-34"} & \texttt{4}\\\bottomrule\end{tabular}\end{center}
\ifpauta\begin{solucionbox}\begin{lstlisting}[style=customc]
int contar_cifras(const char texto[]) {
    int total = 0;
    for (int i = 0; texto[i] != '\0'; i++) {
        if (texto[i] >= '0' && texto[i] <= '9') {
            total++;
        }
    }
    return total;
}
\end{lstlisting}La comparación de rango reconoce cualquier carácter entre cero y nueve en ASCII.\end{solucionbox}\else\vspace{3cm}\fi
\end{enumerate}
\end{document}
